% Figura 3.x · Ponderación de clases frente a submuestreo (Albacete)
% Recuentos exactos de la clase de Albacete (24, 716 y 176 de 916 tramos).
% Pesos relativos: inversamente proporcionales al número de tramos de cada clase.
\begin{figure}[htbp]
  \centering
  \begin{tikzpicture}[
      font=\sffamily\scriptsize,
      x=0.0045cm, % 716 tramos ≈ 3,2 cm
      barra/.style={minimum height=4.2mm, inner sep=0pt, anchor=west},
      etq/.style={anchor=east, text=tinta},
      val/.style={anchor=west, text=tinta-suave},
      tit/.style={anchor=west, font=\sffamily\small\bfseries, text=uoc-azul},
      sub/.style={anchor=north west, text=tinta-suave, align=left},
    ]
    % ---------------- Panel izquierdo: ponderación
    \node[tit] at (0,1.45) {Ponderación de clases};
    \node[sub] at (0,1.2) {Se usan todos los tramos;\\los errores en A pesan más};
    \foreach \cl/\n/\w/\col/\y in {A/24/{alto}/acc-accesible/0, B/716/{bajo}/acc-practicable/-0.55, C/176/{medio}/acc-no-accesible/-1.1}{
      \node[etq] at (-6,\y) {\cl};
      \fill[\col] (0,\y-0.21) rectangle (\n,\y+0.21);
      \node[val] at (\n+8,\y) {\n\ tramos · peso \w};
    }
    \node[sub] at (0,-1.4) {Total: 916 tramos para entrenar};

    % ---------------- Panel derecho: submuestreo
    \begin{scope}[xshift=8.6cm]
      \node[tit] at (0,1.45) {Submuestreo};
      \node[sub] at (0,1.2) {Se reducen B y C hasta\\igualar la clase A};
      \foreach \cl/\n/\col/\y in {A/24/acc-accesible/0, B/716/acc-practicable/-0.55, C/176/acc-no-accesible/-1.1}{
        \node[etq] at (-6,\y) {\cl};
        % parte descartada
        \fill[rejilla] (0,\y-0.21) rectangle (\n,\y+0.21);
        \draw[filete, line width=0.3pt, dashed] (0,\y-0.21) rectangle (\n,\y+0.21);
        % parte que se conserva
        \fill[\col] (0,\y-0.21) rectangle (24,\y+0.21);
      }
      \node[val] at (32,0) {24 tramos};
      \node[val] at (724,-0.55) {24 de 716};
      \node[val] at (184,-1.1) {24 de 176};
      \node[sub] at (0,-1.4) {Total: 72 tramos para entrenar};
      \node[anchor=west, text=tinta-suave] at (0,-2.05) {\muestracolor{rejilla}\ tramos descartados};
    \end{scope}

  \end{tikzpicture}
  \caption{Ponderación de clases frente a submuestreo con los tramos de Albacete.}
  \label{fig:ponderacion-submuestreo}
  \fuente{elaboración propia a partir de la distribución de la figura~\ref{fig:distribucion-clases}, con ayuda de Claude. El peso de cada clase es inversamente proporcional a su número de tramos.}
\end{figure}
